% !TEX program = xelatex
\documentclass[11pt]{article}

% ============ 页面边距 ============
\usepackage[letterpaper, margin=0.6in]{geometry}

% ============ 颜色系统 ============
\usepackage{xcolor}
\definecolor{nameColor}{RGB}{26, 42, 58}       % 深空蓝（姓名）
\definecolor{sectionColor}{RGB}{44, 62, 107}   % 深海蓝（章节标题）
\definecolor{accentColor}{RGB}{58, 110, 165}   % 科技蓝（链接/高亮）
\definecolor{dateColor}{RGB}{74, 90, 106}      % 铁灰（日期/次要信息）
\definecolor{bodyColor}{RGB}{45, 45, 45}       % 深灰（正文）
\definecolor{lightGray}{RGB}{208, 215, 224}    % 浅灰（分割线）

% ============ 字体配置 ============
\usepackage{fontspec}
\usepackage{xeCJK}

% 英文：Source Code Pro（自动匹配系统里存在的变体）
\setmainfont{Source Code Pro}

% 中文：PingFang SC
\setCJKmainfont{PingFang SC}

% 等宽字体
\setmonofont{Source Code Pro}

% ============ 间距与格式 ============
\setlength{\parindent}{0pt}
\setlength{\parskip}{5pt}
\usepackage{enumitem}
\setlist{nosep, left=0pt, itemsep=2pt, label=\textbullet}
\usepackage{url}
\usepackage{hyperref}

% ============ 自定义命令 ============
\newcommand{\sectionheader}[1]{%
    \vspace{10pt}%
    \noindent\textcolor{sectionColor}{\textbf{\Large #1}}%
    \vspace{2pt}%
    \hrulefill
    \vspace{6pt}%
}

\newcommand{\entry}[5]{%
    \noindent\textcolor{nameColor}{\textbf{#1}} \hfill \textcolor{dateColor}{\textit{#2}} \\
    \textcolor{dateColor}{\textit{#3}} \hfill \textcolor{dateColor}{\textit{#4}} \\
    \textcolor{bodyColor}{#5}%
    \vspace{6pt}%
}

% ============ 页眉禁用 ============
\usepackage{fancyhdr}
\pagestyle{empty}

% ============ 文档开始 ============
\begin{document}

% ------------- 个人信息 -------------
\begin{center}
    {\Huge \textcolor{nameColor}{\textbf{张 某}}} \\[6pt]
    \textcolor{dateColor}{
        \href{mailto:zhangmou@email.com}{zhangmou@email.com} \quad | \quad
        +86 138 0000 0000 \quad | \quad
        \href{https://github.com/zhangmou}{GitHub: zhangmou}
    } \\[2pt]
    \textcolor{accentColor}{\textbf{端侧AI部署工程师 · 北京}}
\end{center}

% ------------- 教育经历 -------------
\sectionheader{教育经历}
\entry{北京大学}{2016.09 -- 2020.06}
    {计算机科学与技术 · 学士}{北京}
    {主修：计算机组成原理、操作系统、数据结构与算法、机器学习}

% ------------- 工作经历 -------------
\sectionheader{工作经历}

\entry{某互联网科技公司}{2022.07 -- 至今}
    {高级前端 / AI应用工程师}{北京}
    {• 主导 React Native + TensorFlow.js 端侧AI应用架构设计，实现视频流实时视觉重建与3D渲染交互。 \\
     • 优化移动端内存管理，通过张量生命周期控制和帧率降采样，解决OOM崩溃问题，推理帧率稳定在30fps。 \\
     • 设计并实现模型动态加载方案，将安装包体积压缩至原体积的30\%，用户下载耗时降低60\%。}

\entry{某创业公司}{2020.07 -- 2022.06}
    {前端开发工程师}{上海}
    {• 负责公司核心产品 Web 端与 React Native App 开发。 \\
     • 实现基于 WebSocket 的实时数据看板，支持10万级QPS数据刷新。}

% ------------- 项目经历 -------------
\sectionheader{项目经历}

\entry{视觉重建3D App (React Native + Expo + TF.js)}{2025.12 -- 2026.06}
    {个人项目 \quad | \quad 已上架App Store \& Google Play}{}
    {• 基于 TensorFlow.js 实现手机摄像头视频流的实时语义分割与3D场景重建。 \\
     • 集成 Three.js 渲染引擎，实现真实视频与抽象3D模型的上下分屏实时同步显示。 \\
     • 攻克 iPhone 端推理内存瓶颈，利用 tf.tidy() 与 requestAnimationFrame 循环实现稳定运行。}

\entry{端侧推理框架性能对比工具}{2025.08 -- 2025.11}
    {个人项目}{}
    {• 开发跨平台性能测试工具，对比 TFLite / NCNN / MNN 在移动端的推理延迟与功耗。 \\
     • 使用 Python 自动化收集数据，生成可视化报告，为技术选型提供依据。}

% ------------- 专业技能 -------------
\sectionheader{专业技能}
\begin{itemize}
    \item \textcolor{nameColor}{\textbf{编程语言：}} \textcolor{bodyColor}{JavaScript / TypeScript，Python，C++（基础阅读与调试）}
    \item \textcolor{nameColor}{\textbf{移动端 \& 前端：}} \textcolor{bodyColor}{React Native，Expo，React}
    \item \textcolor{nameColor}{\textbf{AI / 端侧推理：}} \textcolor{bodyColor}{TensorFlow.js，TensorFlow Lite，WebGL 后端，模型量化与压缩}
    \item \textcolor{nameColor}{\textbf{图形学：}} \textcolor{bodyColor}{Three.js，WebGL，3D场景构建与渲染管线基础}
    \item \textcolor{nameColor}{\textbf{其他：}} \textcolor{bodyColor}{Git，Xcode，Android Studio，Metro}
\end{itemize}

% ------------- 其他 -------------
\sectionheader{其他}
\begin{itemize}
    \item \textcolor{bodyColor}{2024年 公司年度技术创新奖（端侧AI落地项目）}
    \item \textcolor{bodyColor}{持有软件设计师中级证书}
\end{itemize}

\end{document}